% EN49 c03 | 原书 原书 p.22–28（PDF 33–39） | 第 1 章 §1.6–1.7
% 骨架文件：标题行已设定，正文由译者填写。每页起首用 % ===== 原书 p.N ===== 注释标记。
\subsection{函子式镜像对称的结果}

% ===== 原书 p.22 =====
我们的第二组结果，讨论镜像对称 A 侧上各自然函子的函子式镜像对称。这些函子由 Auroux 在 [Aur18] 中草图性地提出，把 $\cW((\C^*)^n\setminus H)$ 与 $\cW(H)$、$\cW((\C^*)^n,f)$ 联系起来：
\begin{itemize}
\item （$\Z/2$-分次）\emph{限制函子} $\rho\colon\cW((\C^*)^n\setminus H)\to\cW(H)$，它把 $(\C^*)^n\setminus H$ 中的拉格朗日量送到该拉格朗日量沿被去掉的纤维 $H$ 的“端”（\ref{def:8-2-3}）；
\item \emph{提升函子} $j\colon\cW(H)\to\cW((\C^*)^n\setminus H)$：沿一条从 $0$ 到 $\infty$、避开 $f$ 全部临界值的射线，把 $H$ 中的拉格朗日量平行移动（\ref{def:9-0-1}）；
\item \emph{无穷远处缠绕函子} $\alpha_\infty\colon\cW((\C^*)^n,f)\to\cW((\C^*)^n\setminus H)$，它开启绕无穷远的缠绕（\ref{def:7-1-3}, (1)）；
\item \emph{绕 $H$ 的缠绕函子} $\alpha_0\colon\cW((\C^*)^n,f)\to\cW((\C^*)^n\setminus H)$，它开启绕被去掉的纤维 $H$ 的缠绕（\ref{def:7-1-3}, (2)）。
\end{itemize}
下文各节将详细构造这些函子的各个版本。Auroux 猜想这些函子应满足关系
\begin{enumerate}
\item $\rho\circ\alpha_0\simeq\cap$；
\item $\rho\circ\alpha_\infty\simeq 0$；
\item $\rho\circ j\simeq\id$；
\end{enumerate}
以及正合三角形

\transnote{背景直觉：四个函子里，$\rho$ 与 $j$ 是“看端”与“平行移动”这类直观操作；$\alpha_0$ 与 $\alpha_\infty$ 则是把 $\cW((\C^*)^n,f)$ 中已有的两类缠绕“打开”的函子。三条关系式的含义是：先开启缠绕再看端，等于 $\cap$、等于零、或等于恒等。}

\transnote{易错预警：$j\cap$ 是函子 $j$ 与 cap 函子 $\cap$ 的复合，不是集合的交。$\cap$ 是 Abouzaid–Ganatra 的 cap 函子，与 §1.5.1 中交换图表 (1.2) 里的 $\cup$ 相对；两者分别对应几何中的 cap 与 cup。}

% ===== 原书 p.23 =====
\[
\begin{tikzcd}[row sep=1.0em, column sep=3.6em]
{j\cap} \arrow[rr,"+1"] & & \alpha_\infty \arrow[dl] \\
& \alpha_0 \arrow[ul] &
\end{tikzcd}
\]

\transnote{含糊点拨：原文只把三角形画出来、未写循环方向。按箭头读，它是 $\alpha_0\to j\cap\xrightarrow{+1}\alpha_\infty\to\alpha_0$：$+1$ 标在从 $j\cap$ 出发的那条箭头上，表示三角的第三个映射带分次移位 $1$。}

Auroux 随后就上述函子及其在代数几何中的镜像提出四条猜想。我们在此复述，并指出本论文讨论各条猜想的章节。

\begin{conjecture}[{[Aur18]}]\label{conj:aur-funct}
以下函子在镜像对称下相对应：
\begin{enumerate}
\item （Conjecture 1.2）函子 $\rho$ 镜像于商函子 $q\colon\operatorname{Coh}(Z)\to\operatorname{Sing}(Z)$ 与 Kn{\"o}rrer 周期性等价 $\operatorname{Sing}(Z)\to\operatorname{Coh}(D)$ 的复合（§8）；
\item （Conjecture 1.3 (1)）函子 $\alpha_0$ 镜像于 $i_{X*}$，其中 $i_X\colon X\to Z$ 是包含（§7）；
\item （Conjecture 1.3 (2)）函子 $\alpha_\infty$ 镜像于 $\pi_X^*$，其中 $\pi_X\colon Z\to X$ 是投影（§8）；
\item （Conjecture, §6.2）函子 $j$ 镜像于 $j_{D*}p_D^*(K_X^{-1}|_D\otimes\cdot\,)$，其中 $j_D\colon K_X|_D\to Z$、$p_D\colon K_X|_D\to D$ 分别是包含与投影（§9）。
\end{enumerate}
\end{conjecture}

\transnote{易错预警：条目开头的“（Conjecture 1.2）”“（Conjecture 1.3 (1)）”是 Auroux 论文 [Aur18] 内部的猜想编号，不是本书编号。译文照字面保留，不要误加 \texttt{\textbackslash ref}。}

本论文在按分次数据作适当修改之后证明这些猜想。本论文第二部分的主要结果可概括如下：

\begin{thmstar}
在 \ref{thm:7-2-4} 所述 [Kuw20] 与 [GS22] 的镜像对称等价，与 \ref{thm:6-1-16} 及 \ref{lem:6-1-7} 的 Kn{\"o}rrer 周期性等价复合之后：
\begin{enumerate}
\item （\ref{thm:7-2-4}, (2)）函子 $\alpha_0\colon\cW((\C^*)^n,f)\to\cW_\infty((\C^*)^n\setminus H)$ 镜像于 $i_{X*}$，其中 $i_X\colon X\to Z$ 是包含；
\item （\ref{thm:8-2-1}）函子 $\alpha_\infty\colon\cW((\C^*)^n,f)\to\cW_0((\C^*)^n\setminus H)$ 镜像于 $\pi_X^*[1]$，其中 $\pi_X\colon Z\to X$ 是投影；
\item （\ref{thm:8-2-5}）函子 $\rho\colon\overline{\cW}_0((\C^*)^n\setminus H)\to\overline{\cW}(H)$ 镜像于商函子 $q\colon\overline{\operatorname{Coh}}(Z)\to\overline{\operatorname{Sing}}(Z)$；
\item （\ref{thm:8-3-6}）$\rho$ 的分次提升 $\cW_{0,\infty}((\C^*)^n\setminus H)\to\overline{\cW}(H)$，是 $q\colon\overline{\operatorname{Coh}}_{\mathbb G_m}(Z)\to\overline{\operatorname{Sing}}_{\mathbb G_m}(Z)$ 的分次提升的（分次）镜像；
\item （\ref{thm:9-0-2}）函子 $j$ 镜像于 $j_{D*}p_D^*$，其中 $j_D\colon K_X|_D\to Z$、$p_D\colon K_X|_D\to D$ 分别是包含与投影。
\end{enumerate}
此外还有正合三角形（\ref{cor:9-0-4}）：

% ===== 原书 p.24 =====
\[
\begin{tikzcd}[row sep=1.0em, column sep=3.6em]
{j\cap} \arrow[rr,"+1"] & & \alpha_0 \arrow[dl] \\
& \alpha_\infty \arrow[ul] &
\end{tikzcd}
\]
\end{thmstar}

\transnote{易错预警：定理中 $\cW_\infty$ 与 $\cW_0$ 的下标，对应 §1.5 定义的两个分次 $\eta_\infty$ 与 $\eta_0$；它们与函子名 $\alpha_\infty$、$\alpha_0$ 的下标一致，但两者是不同层面的对象，不可混同。}

\transnote{含糊点拨：$\overline{\cW}$、$\overline{\operatorname{Coh}}$ 上的横杠，是 §1.5 中由 $A_\infty$ 分次模 $2$ 约化得到的记号，表示微分 $\Z/2$-分次范畴，并非拓扑闭包。}

\begin{remarkx}\label{rem:aur-monodromy}
我们的函子与 Auroux 在 [Aur18] 中所定义的函子，相差一次单值变换的作用：Auroux 定义 $\alpha_0,\alpha_\infty$ 时，要求拉格朗日量都先在底空间沿负实无穷方向作平行移动；我们改以扇形包含来定义 $\alpha_0,\alpha_\infty$，于是拉格朗日量分别直接走向负实无穷与正实无穷（参 [Aur18, Fig. 4]）。这意味着我们的结果有时与 [Aur18] 中的猜想相差一次单值变换的作用（参 \ref{thm:9-0-2}）与一次分次平移。
\end{remarkx}

\subsubsection{函子式镜像对称的证明技巧}

% ===== 原书 p.25 =====
本节稍作启发式地综述本论文第二部分用到的一些主要技巧。第一个技巧是 [GPS19, Theorem 1.28] 中 wrapped Fukaya 范畴的同伦推出（homotopy pushout）定理，以及 [GR17, Proposition 4.7.2.2] 中相干层范畴的真下降（proper descent）定理。

$(\C^*)^n\setminus H$ 有两种表示：$\tilde f_0=0$ 或 $\tilde f_\infty=0$。二者自然给出 $\C$ 上两个（奇异）辛纤维丛，分别由 $f$ 与 $1/f$ 给出；同时也给出 $(\C^*)^n\setminus H$ 作为 Liouville 扇形粘合的两个表示。这两个表示又分别对应两个镜像空间 $\tilde Z$ 与 $Z$ 分解为不可约分支。然而，这种辛粘合难以实现。

\transnote{背景直觉：同伦推出定理在 A 侧负责“拼接”，真下降定理在 B 侧负责“拼接”；两侧的拼法在镜像对称下互相对应，这正是本节全部技巧的骨架。}

\transnote{易错预警：这里的“两个（奇异）辛纤维丛”不是同一个纤维化的两种写法，而是 $f$ 与 $1/f$ 两个不同的纤维化；它们的奇异纤维位置互为颠倒，这正是下文“交换零与无穷”的来源。}

给定仿射簇 $X$ 上的多项式函数 $W\colon X\to\C$ 与 Stein 函数 $\phi\colon X\to\R$，实超曲面 $F=\{\operatorname{Re}(W)=C\}$ 应当是扇形超曲面的原型例子：它把 $X$ 表示成两个扇形 $\{\operatorname{Re}(W)\le C\}$ 与 $\{\operatorname{Re}(W)\ge C\}$ 沿 $F$ 的扇形粘合。借助函数 $I=\operatorname{Im}(W)$，$F$ 满足 [GPS20a, Definition 1.2] 的第二个条件；但该定义隐含要求 $F$ 在无穷远处切于 $X$ 的 Liouville 向量场。这一条件一般无法保证，即便把 $X$ 上的 Liouville 结构作形变也未必可行。存在这种简单 Liouville 形变的判据，将出现在第二作者及合作者的未来工作中。这里我们描述另一种构造，它给出相同的结果（§7）。

\transnote{含糊点拨：“$F$ 满足定义的第二个条件”并不等于“$F$ 是扇形超曲面”：缺的是切于 Liouville 向量场这一条。原文用 while 转折，意在强调这个隐含条件才是障碍。}

接下来，我们不必沿整个超曲面 $F$ 修正 Liouville 结构，而可以应用 Sylvan 的一个结果 [Syl19b, Prop 2.6]，在无穷远处某条纤维的邻域内给出标准 Liouville 形式。

% ===== 原书 p.26 =====
令 $H=W^{-1}(R_1)\cap\{\phi\le R_2-\eps\}$，则该命题给出：对某个 $\rho>0$（且 $R_1,R_2$ 充分大），$H$ 有一个标准 Liouville 邻域
\[
U=\left(\hat H\times\{z\in\C:\operatorname{Re}(z)>-\rho\},\ \lambda_X|_H+\lambda_\C^{\operatorname{std}}\right),
\]
它在完备化 $\hat X$ 中真嵌入（见\ref{fig:7-1}）。须提醒读者：投影 $z\colon U\to\C$ 未必与原映射 $W\colon X\to\C$ 相容，尽管 $z$ 的纤维将与 $W$ 的纤维同痕。

\transnote{易错预警：$H=W^{-1}(R_1)\cap\{\phi\le R_2-\eps\}$ 里的 $R_1$ 与 $R_2$ 是两个不同的大参数，分别控制纤维与截断，不能混同；$\rho$ 也与限制函子 $\rho$ 同名，但此处是正实数。}

\transnote{背景直觉：Sylvan 的命题只保证“无穷远处某条纤维附近”有标准 Liouville 形式，所以 $z$ 只能与原 $W$ 的纤维同痕、而不与之相等。这正是这种构造的代价，也是它的用处。}

于是我们得到 $X$ 内真嵌入的扇形超曲面 $\{\operatorname{Re}(z)=0\}\subseteq U$。在我们的情形中取 $W=f$，这就把 $(\C^*)^n$ 表示成 $((\C^*)^n,f)$ 与 $(H\times\C,z)$ 的平凡粘合。若改在邻域 $U$ 内去掉超曲面 $H$，则 $(\C^*)^n\setminus H$ 可表示成 $((\C^*)^n,f)$ 与 $(H\times\C^*,z)$ 的扇形粘合（\ref{prop:7-1-1}）。

同一论证对函数 $1/W\colon X\to\C$ 不成立，因为 $1/W$ 在 $X$ 上并非处处有定义，也不再是多项式。我们可以换一种做法：在 $W$ 诱导的 $X$ 的扇形分解中“交换零与无穷”，由此构造 $1/W$ 诱导的 $X$ 的扇形分解。更确切地说，这意味着从先前的构造出发，施加由 $f$ 的整体单值性给出的纤维 $H$ 的同痕。应用 [GPS20a] 的唯一性结果可知，总能找到一族“跟随”该同痕的扇形超曲面。所得的扇形分解，把 $(\C^*)^n\setminus H$ 表示成 $((\C^*)^n,f)$ 与 $\C^*$ 上一个 $H$-丛的扇形粘合，该丛由 $f$ 的整体单值性扭成（见\ref{fig:7-6}、\ref{prop:7-1-2}）。

\transnote{易错预警：$1/W$ 的问题不在“值太大”，而在于它在 $X$ 上并非处处有定义、且不再是多项式，故同伦推出定理的前提失效，只能改用单值性扭成的粘合。}

本论文用到的第二个主要技巧，是带规范群的 LG 模型的导出 Kn{\"o}rrer 周期性（参 [Hir17]）。由 $1/f$ 诱导的 $(\C^*)^n\setminus H$ 的扇形粘合，在镜像对称下对应 $Z$ 由其不可约分支的 B 侧粘合（\ref{thm:7-2-4}）；另一方面，非常规的分次 $\eta_0$ 意味着：由 $f$ 诱导的 $(\C^*)^n\setminus H$ 的扇形分解，给出带分次 $\eta_0$ 的 $(\C^*)^n\setminus H$ wrapped Fukaya 范畴与导出概型 $\tilde Z$ 上相干层之间的一个等价。

% ===== 原书 p.27 =====
为把这一结果拉回到 [Aur18] 关于 $Z$ 的原始镜像对称猜想，我们需要 $\operatorname{Coh}(Z)$ 与 $\operatorname{Coh}(\tilde Z)$ 之间的等价。这两个概型由 Koszul 对偶与导出 Kn{\"o}rrer 周期性的组合联系起来：二者有公共的 Kn{\"o}rrer 稳定化，稳定到 $K_X\times\mathbb A^1$ 上 $syz$ 的（分次）矩阵分解范畴，其中 $y,z$ 分别记 $K_X$ 与 $\mathbb A^1$ 的纤维上的坐标（\ref{thm:6-1-16}）。

\transnote{含糊点拨：$\operatorname{Coh}(Z)$ 与 $\operatorname{Coh}(\tilde Z)$ 之间的等价并非恒等，而是经“公共 Kn{\"o}rrer 稳定化”过渡得到的；原文用 stabilization 一词，指两个概型各自再添一对互为对偶的坐标后同构。}

$\operatorname{Coh}(Z)$ 与 $\operatorname{Coh}(\tilde Z)$ 之间的这个等价具有出人意料的性质，使它成为把 $f$ 换成 $1/f$ 的代数类比。该等价把函子 $i_{\tilde X*}$ 送到 $\pi_X^*[1]$（\ref{lem:8-1-2}），反之亦然；因此 A 侧诱导的等价交换函子 $\alpha_0$ 与 $\alpha_\infty$（\ref{thm:8-2-1}），在某种意义上“交换零与无穷”。这类耐人寻味的非几何等价的一般意义，将留待未来工作讨论。

\transnote{背景直觉：“把 $f$ 换成 $1/f$”是 A 侧的整体单值性操作，其在 B 侧的代数类比就是 $\operatorname{Coh}(Z)\leftrightarrow\operatorname{Coh}(\tilde Z)$。会“交换零与无穷”，指两个缠绕函子 $\alpha_0,\alpha_\infty$ 在该等价下互换。}

\subsection{与已发表工作的关系}

本论文的结果取自两篇文章 [Jef22] 与 [GJ22]，后者是与 Benjamin Gammage 的合作工作。下面列出本论文各节的来源：
\begin{itemize}
\item §1 综合了 [Jef22] §1 与 [GJ22] §1 的内容（含与 Benjamin Gammage 的合作工作）；
\item §2 综合了 [Jef22] §2 与 [GJ22] §3.1 的内容（含与 Benjamin Gammage 的合作工作）；
\item §3 对应 [Jef22] 的 §3；
% ===== 原书 p.28 =====
\item §4 对应 [Jef22] 的 §4；
\item §5 对应 [Jef22] 的 §5；
\item §6 对应 [GJ22] 的 §2（与 Benjamin Gammage 的合作工作）；
\item §7 对应 [GJ22] 的 §3（与 Benjamin Gammage 的合作工作）；
\item §8 对应 [GJ22] 的 §4（与 Benjamin Gammage 的合作工作）；
\item §9 对应 [GJ22] 的 §5（与 Benjamin Gammage 的合作工作）；
\item 附录 A 对应 [GJ22] 的附录 A（与 Benjamin Gammage 的合作工作）。
\end{itemize}

\transnote{易错预警：本节两套“§N”不可互换——左边的“§1”指本论文的节，括号前的“§1 of [Jef22]”指 [Jef22] 的节。译文一律照字面保留，读者须按出处区分。}
